%=====================================================================
% ESTILO DA CASA -- TEMPLATE BEAMER (padrao da colecao)
% Arquivo central de identidade visual para todos os materiais.
% Azul-petroleo + ambar | tcolorbox | TikZ/pgfplots/circuitikz
%=====================================================================

% Dependencias minimas do proprio estilo.
\RequirePackage{xcolor}
\RequirePackage{tikz}
\usetikzlibrary{calc}
\RequirePackage{tcolorbox}
\tcbuselibrary{skins,breakable}
\RequirePackage{listings}
\RequirePackage{xparse}
\RequirePackage{etoolbox}

\makeatletter

%====================== PALETA OFICIAL ================================
\definecolor{corPrim}{RGB}{13,48,82}      % azul-petroleo
\definecolor{corAccent}{RGB}{222,138,28}  % ambar
\definecolor{corDef}{RGB}{31,111,168}     % azul medio
\definecolor{corNorma}{RGB}{0,121,107}    % verde-petroleo
\definecolor{corEx}{RGB}{46,125,50}       % verde
\definecolor{corFix}{RGB}{204,102,0}      % laranja
\definecolor{corAt}{RGB}{176,58,46}       % vermelho
\definecolor{codBg}{RGB}{252,250,245}
\definecolor{codKw}{RGB}{13,48,82}
\definecolor{codStr}{RGB}{46,125,50}
\definecolor{codCom}{RGB}{120,130,138}

% Compatibilidade com nomes usados historicamente nos materiais.
\colorlet{Petroleo}{corPrim}
\colorlet{PetroleoEscuro}{corPrim!82!black}
\colorlet{PetroleoMedio}{corDef!82!corPrim}
\colorlet{PetroleoClaro}{corDef!7}
\colorlet{Ambar}{corAccent}
\colorlet{AmbarEscuro}{corAccent!72!black}
\colorlet{AmbarClaro}{corAccent!16}
\colorlet{AzulPetroleo}{corPrim}
\colorlet{AzulPetroleoEscuro}{corPrim!82!black}
\colorlet{AzulPetroleoClaro}{corDef!7}
\colorlet{AzulMedio}{corDef}
\colorlet{AzulClaro}{corDef!7}
\colorlet{CinzaEscuro}{black!82}
\colorlet{CinzaClaro}{black!5}
\colorlet{CinzaSuave}{black!3}
\colorlet{CinzaMedio}{black!55}
\colorlet{CinzaTexto}{black!82}
\colorlet{CinzaBox}{black!3}
\colorlet{Verde}{corEx}
\colorlet{Vermelho}{corAt}
\colorlet{Turquesa}{corNorma!70!corDef}
\colorlet{Roxo}{corDef!58!corAt}
\colorlet{ifmagreen}{corPrim}
\colorlet{ifmagreen2}{corDef}
\colorlet{ifmalight}{corDef!7}
\colorlet{accent}{corAccent}
\colorlet{warn}{corAt}
\colorlet{ok}{corEx}
\colorlet{bluex}{corDef}
\colorlet{grayx}{black!58}

% Variantes em minusculas presentes em materiais antigos.
\colorlet{azulPetroleo}{corPrim}
\colorlet{azulPetroleoEscuro}{corPrim!82!black}
\colorlet{azulPetroleoClaro}{corDef!7}
\colorlet{azulMedio}{corDef}
\colorlet{azulClaro}{corDef!7}
\colorlet{azulEscuro}{corPrim!82!black}
\colorlet{azul}{corPrim}
\colorlet{azul2}{corDef}
\colorlet{ambar}{corAccent}
\colorlet{ambarClaro}{corAccent!16}
\colorlet{cinza}{black!5}
\colorlet{cinzaClaro}{black!5}
\colorlet{cinzaMedio}{black!55}
\colorlet{cinzaEscuro}{black!82}
\colorlet{cinzaTexto}{black!82}
\colorlet{verde}{corEx}
\colorlet{verdeclaro}{corEx!10}
\colorlet{vermelho}{corAt}
\colorlet{vermelhoclaro}{corAt!9}
\colorlet{laranja}{corFix}
\colorlet{roxo}{corDef!58!corAt}

%====================== TEMA BEAMER ==================================
\usetheme{default}
\usefonttheme{professionalfonts}
\setbeamercolor{normal text}{fg=black!85,bg=white}
\setbeamercolor{structure}{fg=corPrim}
\setbeamercolor{frametitle}{bg=corPrim,fg=white}
\setbeamercolor{footsec}{bg=corPrim,fg=white!92!corPrim}
\setbeamercolor{alerted text}{fg=corAccent!85!black}
\setbeamerfont{frametitle}{series=\bfseries,size=\large}
\setbeamerfont{footline}{size=\tiny}
\setbeamertemplate{navigation symbols}{}
\setbeamertemplate{caption}{\raggedright\insertcaption\par}
\setbeamersize{text margin left=0.55cm,text margin right=0.55cm}

\setbeamertemplate{itemize item}{\raise0.35ex\hbox{\color{corAccent}\rule{0.62ex}{0.62ex}}}
\setbeamertemplate{itemize subitem}{\color{corDef}\raise0.5pt\hbox{--}}
\setbeamertemplate{itemize subsubitem}{\color{corPrim}\tiny$\circ$}
\setbeamertemplate{enumerate item}{\color{corPrim}\textbf{\insertenumlabel.}}
\setbeamertemplate{section in toc}{%
  \leavevmode\textcolor{corAccent}{\textbf{\inserttocsectionnumber.}}~\inserttocsection\par\vskip1.5pt}
\setbeamertemplate{subsection in toc}{%
  \leavevmode\leftskip=0.95cm\textcolor{corDef}{\textbullet}~\inserttocsubsection\par}

%====================== TITULO DO SLIDE ==============================
\setbeamertemplate{frametitle}{%
  \nointerlineskip%
  \begin{beamercolorbox}[wd=\paperwidth,ht=2.7ex,dp=1.3ex,leftskip=0.55cm,rightskip=0.55cm]{frametitle}%
    \usebeamerfont{frametitle}\insertframetitle%
  \end{beamercolorbox}%
  {\color{corAccent}\rule{\paperwidth}{1.3pt}}\par\vskip2pt%
}

%====================== RODAPE DUPLO ================================
\newcommand{\CasaRodapeEsq}{%
  \ifcsname Disciplina\endcsname\Disciplina\else\insertshorttitle\fi}
\newcommand{\CasaRodapeDir}{%
  \ifcsname TituloCurto\endcsname\TituloCurto\else\insertshorttitle\fi}
\newcommand{\CasaMetaCapa}{%
  \ifcsname Disciplina\endcsname
    \if\relax\detokenize\expandafter{\Disciplina}\relax
    \else
      \Disciplina\ \textbullet\ \insertdate
    \fi
  \else
    \insertdate
  \fi}
\setbeamertemplate{footline}{%
  \leavevmode%
  \hbox{%
    \begin{beamercolorbox}[wd=.5\paperwidth,ht=2.4ex,dp=1ex,leftskip=0.5cm]{footsec}%
      \usebeamerfont{footline}\CasaRodapeEsq%
    \end{beamercolorbox}%
    \begin{beamercolorbox}[wd=.5\paperwidth,ht=2.4ex,dp=1ex,rightskip=0.5cm]{footsec}%
      \hfill\usebeamerfont{footline}\CasaRodapeDir\quad\textbar\quad\insertframenumber%
    \end{beamercolorbox}%
  }\vskip0pt%
}

%====================== DIVISORIAS DE SECAO =========================
\newcommand{\CasaSecNum}{\ifnum\value{section}<10 0\fi\arabic{section}}
\newcommand{\CasaDivisoriaManual}[2]{%
  \begin{frame}[plain,noframenumbering]
    \makebox[0pt][l]{%
      \raisebox{0pt}[0pt][0pt]{%
        \begin{tikzpicture}[remember picture,overlay]
          \fill[corPrim] (current page.south west) rectangle (current page.north east);
          \if\relax\detokenize{#1}\relax\else
            \node[anchor=west,text=corAccent,font=\bfseries\fontsize{40}{40}\selectfont]
              at ([xshift=1.6cm,yshift=1.1cm]current page.west) {#1};
          \fi
          \fill[corAccent] ([xshift=1.7cm,yshift=0.35cm]current page.west)
            rectangle ++(2.4,0.09);
          \node[anchor=west,text=white,font=\bfseries\LARGE,text width=0.75\paperwidth]
            at ([xshift=1.6cm,yshift=-0.5cm]current page.west) {#2};
        \end{tikzpicture}%
      }%
    }%
  \end{frame}%
}
\AtBeginSection[]{\CasaDivisoriaManual{\CasaSecNum}{\insertsection}}

% Alguns materiais antigos usam \secslide{titulo}; outros, \secslide{numero}{titulo}.
\@ifundefined{secslide}
  {\NewDocumentCommand{\secslide}{m g}{\IfNoValueTF{#2}{\CasaDivisoriaManual{}{#1}}{\CasaDivisoriaManual{#1}{#2}}}}
  {\RenewDocumentCommand{\secslide}{m g}{\IfNoValueTF{#2}{\CasaDivisoriaManual{}{#1}}{\CasaDivisoriaManual{#1}{#2}}}}
\@ifundefined{Divisoria}
  {\NewDocumentCommand{\Divisoria}{m g}{\IfNoValueTF{#2}{\CasaDivisoriaManual{}{#1}}{\CasaDivisoriaManual{#1}{#2}}}}
  {\RenewDocumentCommand{\Divisoria}{m g}{\IfNoValueTF{#2}{\CasaDivisoriaManual{}{#1}}{\CasaDivisoriaManual{#1}{#2}}}}

%====================== CAPA ========================================
\defbeamertemplate*{title page}{estilo da casa}{%
  \vspace{0.8cm}%
  {\color{white}\bfseries\huge\inserttitle\par}%
  \vspace{2mm}%
  {\color{corAccent}\rule{3.2cm}{2.2pt}\par}%
  \vspace{4mm}%
  {\color{white!85}\large\insertsubtitle\par}%
  \vspace{8mm}%
  {\color{white}\normalsize\insertauthor\par}%
  \vspace{1mm}%
  {\color{white!70}\small\CasaMetaCapa\par}%
}
\setbeamertemplate{title page}[estilo da casa]

%====================== CAIXAS =======================================
% Materiais antigos usam convencoes diferentes para o argumento opcional
% das caixas. Para preservar a semantica, caixas existentes nao sao
% redefinidas; apenas sua paleta historica e recolorida acima.
\tcbset{
  caixabase/.style={
    enhanced,sharp corners=downhill,boxrule=0.6pt,left=2mm,right=2mm,
    top=1.2mm,bottom=1.2mm,fonttitle=\bfseries\small,
    attach boxed title to top left={xshift=3mm,yshift=-2.4mm},
    boxed title style={sharp corners,boxrule=0pt},
    before skip=3pt,after skip=3pt
  }
}
\@ifundefined{caixaDef}{\newtcolorbox{caixaDef}[1][]{caixabase,colback=corDef!5,colframe=corDef,coltitle=white,colbacktitle=corDef,title=Definição,#1}}{}
\@ifundefined{caixaNorma}{\newtcolorbox{caixaNorma}[1][]{caixabase,colback=corNorma!5,colframe=corNorma,coltitle=white,colbacktitle=corNorma,title=Norma / Método,#1}}{}
\@ifundefined{caixaEx}{\newtcolorbox{caixaEx}[1][]{caixabase,colback=corEx!5,colframe=corEx,coltitle=white,colbacktitle=corEx,title=Exemplo resolvido,#1}}{}
\@ifundefined{caixaFix}{\newtcolorbox{caixaFix}[1][]{caixabase,colback=corFix!5,colframe=corFix,coltitle=white,colbacktitle=corFix,title=Fixação,#1}}{}
\@ifundefined{caixaAt}{\newtcolorbox{caixaAt}[1][]{caixabase,colback=corAt!5,colframe=corAt,coltitle=white,colbacktitle=corAt,title=Atenção,#1}}{}
\@ifundefined{caixaForm}{\newtcolorbox{caixaForm}[1][]{caixabase,colback=corPrim!4,colframe=corPrim,coltitle=white,colbacktitle=corPrim,title=Formulário,#1}}{}
\@ifundefined{caixaObs}{\newtcolorbox{caixaObs}[1][]{caixabase,colback=black!3,colframe=corPrim,coltitle=white,colbacktitle=corPrim,title=Observação,#1}}{}
\@ifundefined{caixaAlerta}{\newtcolorbox{caixaAlerta}[1][]{caixabase,colback=corAt!5,colframe=corAt,coltitle=white,colbacktitle=corAt,title=Atenção,#1}}{}
\@ifundefined{caixaProjeto}{\newtcolorbox{caixaProjeto}[1][]{caixabase,colback=corNorma!5,colframe=corNorma,coltitle=white,colbacktitle=corNorma,title=Projeto prático,#1}}{}
\@ifundefined{caixaProj}{\newtcolorbox{caixaProj}[1][]{caixabase,colback=corNorma!5,colframe=corNorma,coltitle=white,colbacktitle=corNorma,title=Projeto prático,#1}}{}

%====================== CODIGO =======================================
\lstdefinestyle{codCasa}{
  basicstyle=\ttfamily\scriptsize,
  keywordstyle=\color{codKw}\bfseries,
  stringstyle=\color{codStr},commentstyle=\color{codCom}\itshape,
  numbers=left,numberstyle=\tiny\color{codCom},numbersep=5pt,
  frame=single,rulecolor=\color{corPrim!25},backgroundcolor=\color{codBg},
  breaklines=true,showstringspaces=false,tabsize=2,columns=fullflexible
}
\lstdefinestyle{codC}{style=codCasa,language=C}
\lstdefinestyle{codPy}{style=codCasa,language=Python}
\lstdefinestyle{py}{style=codCasa,language=Python}

%====================== EXTENSOES MODULARES ==========================
% Se existir <jobname>_antes_laboratorios.tex, o arquivo e carregado uma
% unica vez imediatamente antes da secao de laboratorios. Isso permite
% ampliar um material sem duplicar o arquivo principal.
\newif\ifCasaExtensaoAntesLabsInserida
\CasaExtensaoAntesLabsInseridafalse
\AtBeginDocument{%
  \IfFileExists{\jobname_antes_laboratorios.tex}{%
    \let\CasaSectionOriginal\section
    \renewcommand{\section}[1]{%
      \ifCasaExtensaoAntesLabsInserida\else
        \ifstrequal{##1}{Laboratórios, projetos e comissionamento}{%
          \global\CasaExtensaoAntesLabsInseridatrue
          \input{\jobname_antes_laboratorios.tex}%
        }{}%
      \fi
      \CasaSectionOriginal{##1}%
    }%
  }{}%
}

\makeatother
